\section{Ordinal norms and effective collapse evaluation}
\label{f:sec:foundations}

The numerical assignment uses ordinals below $\varepsilon_0$, represented
by hereditary Cantor normal forms.  The ordinal $\varepsilon_0$ is the
least fixed point of $\alpha\mapsto\omega^\alpha$.  We distinguish
ordinary ordinal addition and multiplication from their natural
counterparts $\oplus$ and $\otimes$.  Natural addition collects the
Cantor monomials of its operands; natural multiplication distributes over
those monomials and naturally adds their exponents.  In particular both
natural operations are commutative.  All finite tails displayed as
$\lambda+n$ use ordinary addition.  If $\lambda$ is a limit and $n$ is
finite, then $\lambda\oplus n=\lambda+n$.

\begin{definition}[Additive hereditary norm]\label{f:def:additive-norm}
For decreasing exponents and positive integer coefficients put
\[
 N(0)=0,\qquad
 N\left(\sum_{i<r}\omega^{\alpha_i}c_i\right)
       =\sum_{i<r}c_i\bigl(1+N(\alpha_i)\bigr).
\]
Thus $N(n)=n$, $N(\omega)=2$ and $N(\omega^2)=3$.
Coefficients are charged with their full multiplicity.
\end{definition}

\begin{lemma}[Effective ordinal operations and norm bounds]
\label{f:lem:norm-arithmetic}
Every norm ball is finite and effectively enumerable.  Cantor-normal-form
comparison, $\oplus$, $\otimes$, and ordinal exponentiation $2^\alpha$
are effective.  Moreover
\begin{align*}
 N(\alpha\oplus\beta)&=N(\alpha)+N(\beta),\\
 N(\alpha)+N(\beta)-1
 &\le N(\alpha\otimes\beta)\le N(\alpha)N(\beta)
       &&(\alpha,\beta>0),\\
 N(\alpha)&\le2N(2^\alpha),\qquad
 N(2^\alpha)\le2^{N(\alpha)}.
\end{align*}
\end{lemma}
\begin{proof}
An exponent in a nonzero ordinal of norm at most $b$ has norm at most
$b-1$.  Induction on $b$ therefore enumerates the finite set of possible
exponents and the finite lists of coefficients whose total charge is at
most $b$.  Recursive lexicographic comparison decides their order.
The definitions of natural sum and product give finite algorithms.

For the norm bounds, expand coefficients into repeated monomials.
Suppose the two operands have respectively $r,s$ monomials, and let
$A,B$ be the sums of the norms of their exponents.  Their norms are
$r+A$ and $s+B$, while the product has norm $rs+sA+rB$.
Collecting equal monomials preserves these totals.  The upper product
bound follows by expanding $(r+A)(s+B)$; for the lower bound the
difference is
\[
 (r-1)(s-1)+(s-1)A+(r-1)B\ge0.
\]
Norm additivity for natural sum follows in the same way.

Write $\alpha=\omega\xi+n$, $n<\omega$, using ordinary
multiplication.  Then
\[
 2^\alpha=\omega^\xi 2^n,\qquad
 N(2^\alpha)=2^n(1+N(\xi)).
\]
Multiplication on the left by $\omega$ changes the exponent of each
monomial from $\gamma$ to $1+\gamma$.  Its norm stays the same or
increases by one, so
$N(\xi)\le N(\omega\xi)\le2N(\xi)$.
The last two asserted bounds now follow from
$1+a\le2^a$ for $a\in\N$ and
$2N(\xi)+n\le2^{n+1}(1+N(\xi))$.
These are the norm estimates of \cite[Proposition~2.2]{ww}, with the
additive convention made explicit.
\end{proof}

For reference, the finite hierarchy used throughout is
\[
 F_0(n)=2^n,\qquad F_{j+1}(n)=F_j^{\,n+1}(n),\qquad
 G(n)=F_5(n+100).
\]
We fix $f=G$ in this paper.  No residue correction or other change to the
control is made.  The estimates used below are proved in
Section~\ref{m:sec:control}.  Superscripts on a function denote finite
iteration, with the zeroth iterate equal to the identity.

\begin{definition}[Norm-controlled collapse]\label{f:def:collapse}
For $\alpha<\varepsilon_0$ set
\[
 \psi(\alpha)=\max\bigl(\{0\}\cup
   \{\psi(\beta)+1:\beta<\alpha,
                    \ N(\beta)\le f(N(\alpha))\}\bigr).
\]
Ordinal recursion defines this function, since each displayed maximum
is finite.
\end{definition}

\begin{lemma}[Basic collapse estimates]\label{f:lem:collapse-estimates}
If $\beta<\alpha$ and $N(\beta)\le f(N(\alpha))$, then
$\psi(\beta)+1\le\psi(\alpha)$.  For $n,k\in\N$,
\[
 \psi(n)=n,\qquad N(\alpha)\le\psi(\alpha),\qquad
 \psi(\alpha)+k\le\psi(\alpha+k).
\]
If $\alpha\ge\omega$, then $\psi(\alpha)>f(N(\alpha))$.
\end{lemma}
\begin{proof}
The first assertion is a term of the defining maximum.  Induction on
$n$ gives the finite values: every predecessor of $n$ is finite and
its norm is below $n\le f(n)$.  If $\alpha\ge\omega$, the finite
ordinal $f(N(\alpha))$ is an admitted predecessor, which proves the
strict estimate and hence the norm bound.  The finite case of the norm
bound has already been proved.  Finally $N(\alpha+k)=N(\alpha)+k$,
so $\alpha+k$ is an admitted predecessor of $\alpha+k+1$.
Induction on $k$ proves the finite-tail inequality.
\end{proof}

In particular a fixed-head collapse is increasing in its finite tail.
The lemma does not assert monotonicity of $\psi$ on arbitrary ordinal
arguments.  The norm premise must be verified whenever an ordinal
inequality is passed through the collapse.

\subsection{An exact maximum-predecessor algorithm}

For $\alpha>0$ and $b\in\N$ define
\[
 M(\alpha,b)=\max\{\beta<\alpha:N(\beta)\le b\},\qquad
 L(e,b)=M(\omega^e,b).
\]
The defining set for $M$ contains zero and is finite.  We also put
$M(0,b)=0$ as a convention; it is not a predecessor assertion at zero.

\begin{theorem}[Greedy maximum]\label{f:thm:maximum-predecessor}
The following simultaneous recursion computes $M$ and $L$ exactly.
Write $\alpha=\omega^e c+r$, where $c>0$ and $r<\omega^e$, and
put $w=1+N(e)$.
\begin{enumerate}[label=(\roman*)]
\item $M(\alpha,0)=0$ and $L(0,b)=L(e,0)=0$.
\item If $r>0$ and $cw\le b$, then
\[
 M(\alpha,b)=\omega^e c+M(r,b-cw).
\]
\item Otherwise set
\[
 d=\begin{cases}
 \min\{c-1,\lfloor b/w\rfloor\},&r=0,\\
 \min\{c,\lfloor b/w\rfloor\},&r>0.
 \end{cases}
\]
Then $M(\alpha,b)=\omega^e d+L(e,b-dw)$, with a zero block
omitted.
\item If $e,b>0$, put $a=M(e,b-1)$, $v=1+N(a)$ and
$(q,s)=\operatorname{divmod}(b,v)$.  Then
\[
 L(e,b)=\omega^a q+L(a,s).
\]
\end{enumerate}
\end{theorem}
\begin{proof}
If an affordable nonzero tail is present, retaining the entire leading
block and decreasing the tail dominates every predecessor with a smaller
leading block.  This proves (ii).  Otherwise the coefficient in (iii)
is the largest affordable coefficient compatible with the strict upper
bound.  Its remaining tail is less than $\omega^e$ and has precisely
the remaining norm budget.

For (iv), a nonzero candidate below $\omega^e$ with budget $b$ has
a leading exponent $a'<e$ satisfying $N(a')\le b-1$.
The largest possible leading exponent is therefore $a=M(e,b-1)$.
Cantor comparison then maximizes its coefficient and the smaller tail
in that order, yielding $q$ and $L(a,s)$.

Termination follows from the lexicographic measure (budget, phase),
with $M$ above $L$ at equal budget.  Clause (ii) lowers the budget;
(iii) either lowers it or passes from $M$ to $L$.  In (iv) the first
recursive call uses $b-1$, while $v\le b$ and $q\ge1$ give $s<b$
for the second.  This proves both effectiveness and exactness.
\end{proof}

\begin{lemma}[Limit saturation and monotonicity]\label{f:lem:limit-saturation}
For every nonzero limit $\lambda$ and every $b\in\N$,
\[
 N(M(\lambda,b))=b.
\]
For fixed $\lambda$ this maximum is strictly increasing in $b$.
For arbitrary bounds the function $M$ is weakly increasing in each
argument wherever the predecessor interpretation is defined.
\end{lemma}
\begin{proof}
If $\eta=M(\lambda,b)$ has unused budget $d=b-N(\eta)>0$,
then $\eta+d<\lambda$ and $N(\eta+d)=b$, contradicting
maximality.  For $b<c$, the predecessor $M(\lambda,b)+1$ is
below $\lambda$ and has norm at most $c$, proving strict increase.
The last assertion follows by inclusion of the defining candidate sets.
\end{proof}

\begin{theorem}[Single-predecessor collapse recurrence]
\label{f:thm:single-predecessor}
For every nonzero limit $\lambda$ and $n\in\N$,
\[
 \psi(\lambda+n)
   =\psi\bigl(M(\lambda,f(N(\lambda)+n))\bigr)+1.
\]
\end{theorem}
\begin{proof}
Induct on $n$.  Put $b=f(N(\lambda)+n)$ and
$\eta=M(\lambda,b)$.  It is an admitted predecessor, of norm $b$,
so it supplies the lower inequality.  Any other admitted $\beta<\lambda$
satisfies $\beta\le\eta$.  If $\beta<\eta$, then
$N(\beta)\le b=N(\eta)\le f(N(\eta))$; controlled strictness
therefore gives $\psi(\beta)+1\le\psi(\eta)$.

The remaining predecessors are $\lambda+j$, $j<n$.  For such a $j$
put $b_j=f(N(\lambda)+j)$ and $\eta_j=M(\lambda,b_j)$.
Strict growth of $f$ and limit saturation give
$\eta_j<\eta$ and $N(\eta_j)=b_j<b$.
The induction hypothesis and controlled strictness imply
\[
 \psi(\lambda+j)=\psi(\eta_j)+1\le\psi(\eta).
\]
Thus every contribution to the original maximum is at most
$\psi(\eta)+1$, which proves the reverse inequality.
\end{proof}

\begin{corollary}[Closed evaluation and its norm orbit]
\label{f:cor:closed-evaluation}
A closed collapse value is computable by the following deterministic
descent.  At a nonfinite state $\alpha=\lambda+n$, replace it by
$M(\lambda,f(N(\alpha)))$ and add one to an event counter.
On reaching a finite state $Z$, return $Z$ plus that counter.
If the initial norm is $t$ and the number of events is $D$, then
\[
 Z=f^D(t),\qquad \psi(\alpha)=Z+D.
\]
For a nonfinite initial state one also has $0<D<Z$.
\end{corollary}
\begin{proof}
Each step strictly lowers the ordinal, so well-foundedness below
$\varepsilon_0$ proves termination.  The preceding theorem preserves
the value after adding the event count, and finite evaluation gives the
last value.  Limit saturation makes the norm of each new state exactly
$f$ of the preceding norm, including the transition to the finite state.
There is no further event count along the finite predecessor chain.
A nonfinite start has $t\ge2$, and $f(n)>n$ gives
$Z\ge t+D>D$.
\end{proof}

\begin{corollary}[Finite coefficient identity]\label{f:cor:finite-coefficient}
For all $k,n\in\N$,
\[
 \psi(\omega k+n)=f^k(2k+n)+k.
\]
\end{corollary}
\begin{proof}
The case $k=0$ is finite evaluation.  If $k>0$ and
$b=f(2k+n)$, the exact maximum is
$M(\omega k,b)=\omega(k-1)+(b-2(k-1))$.
It has total norm $b$.  The single-predecessor recurrence and induction
on $k$ give the formula.
\end{proof}

The evaluator always calls $M$ with the \emph{limit part} of its
current ordinal.  Calling it with $\lambda+n$ instead would retain
finite predecessors that have already been accounted for in the maximum
argument and would not give the stated norm orbit.
